\section{The curve between the quotes}
\label{sec:clkinterp}

We can now read quoted expiries on the right clock.  But a board
quotes a handful of maturities, and half the daily uses of a term
structure need values \emph{between} them: plotting a dense curve,
marking an unlisted maturity, feeding a model that wants variance at
arbitrary dates.  Between the quotes something must be assumed, and
this section shows that the assumption is the clock choice again ---
in the sharpest form the chapter has to offer.

The rule the clock recommends is simple.  Interpolate \emph{total
variance}, not volatility --- variance is the additive quantity, the
thing that accumulates day by day, while a volatility is a ratio and
ratios do not add.  And interpolate it \emph{linearly in $\tau$}:
between two quoted expiries, let $w$ grow at a constant rate per
variance day.  Constant accrual per unit of the market's clock is
exactly the steady-diffusion picture of \cref{sec:clkwalk}, restated
between quotes; the alternative --- linear in $t$, constant accrual
per calendar day --- is the same picture on the wrong clock, and we
are about to measure the difference.

Build the test so that the truth is known.  A synthetic name runs at
a flat $30\%$ clock volatility --- the previous section's name ---
with one event worth \MacClkClockExtraDays\ extra days at day
\MacClkInterpEventDay.  Three expiries are quoted: 7, 30 and 90 days.
The event falls \emph{between} the second and third quotes: no quoted
expiry sits near it on either side.  Generate the three quotes from
the construction ($w=(0.30)^2\tau$ at each), then hand them --- plus
the event calendar --- to the two interpolation rules and ask each
for the whole dense curve.

\Cref{fig:clkinterp}(a) shows both answers as volatility readings
against calendar days.  Both curves pass through the three quoted
dots --- at a quote there is nothing to interpolate, so the rules
cannot disagree there.  Between 7 and 30 days they also agree: no
event inside, so a variance day is a calendar day and the two rules
coincide.  Everything happens between 30 and 90.  The blue curve ---
linear in $\tau$ --- runs flat at $30\%$ to day \MacClkInterpEventDay,
jumps as the event enters the maturity, and decays along the hump
shape of \cref{sec:clkcrush}: it reproduces the generating truth
exactly (to \MacClkInterpExactBp\ vol bp --- exactness is by
construction here, since the quotes sit on a curve that is itself
linear in $\tau$; the honest general statement closes the section).
The orange dashed curve --- linear in $t$ --- takes the interval's
total variance budget, which is correct at both endpoints, and
spreads it \emph{evenly across all sixty calendar days}.  The event's
four days of variance are smeared over the whole bracket.

Where does the smear hurt, and by how much?  Work the two sides of
the event date; the reading at maturity $T$ (in days) is
$30\%\times\sqrt{x(T)/T}$ with $x(T)$ the interpolated variance-day
count.  First the bracket's budget.  By the clock
\eqref{eq:clkclock}, the 30-day quote carries $x(30)=30$ variance
days and the 90-day quote carries $x(90)=90+4=94$ --- the event is
inside it --- so the bracket holds $94-30=64$ variance days over
$60$ calendar days, and the calendar rule pays them out evenly:
$\tfrac{64}{60}$ variance days per calendar day.  Just \emph{before} day 45, the truth carries no event:
$x_{\rm true}=45$.  The calendar rule has been paying the event out
steadily since day 30, so it carries
$x_{\rm cal}=30+15\cdot\tfrac{64}{60}=46$ --- one phantom variance
day --- and reads
$30\%\sqrt{46/45}=30.33\%$: \MacClkInterpOverBp\ vol bp of phantom
volatility at a maturity that contains no event at all.  Just
\emph{after} day 45, the truth jumps to $x_{\rm true}=49$ while the
calendar rule still shows $x_{\rm cal}=46$, three days short, and
reads $30\%\sqrt{46/45}$ against a truth of $30\%\sqrt{49/45}$:
\MacClkInterpUnderBp\ vol bp \emph{too low} on the first maturities
that do contain the event.  Panel~(b) plots this error across the
whole axis: zero wherever a quote or an eventless bracket pins the
rules together, a sawtooth around the event --- climbing to
$+\MacClkInterpOverBp$, snapping to $-\MacClkInterpUnderBp$, decaying
back to zero at the 90-day quote.

\begin{figure}[htbp]
\centering
\paperfig{\textwidth}{fig_clk_interp.pdf}
\caption{Interpolate in the clock, not the calendar (three quotes,
flat $30\%$ clock truth, event of \MacClkClockExtraDays\ extra days at
day \MacClkInterpOverDay).  (a)~Linear-in-$\tau$ (blue) reproduces the
truth exactly; linear-in-$t$ (orange, dashed) smears the event across
the whole 30--90-day bracket.  Both pass through every quote.
(b)~The calendar rule's reading error: $+\MacClkInterpOverBp$ vol bp
of phantom volatility just before the event date,
$-\MacClkInterpUnderBp$ vol bp just past it.}
\label{fig:clkinterp}
\end{figure}

Both signs are money.  A desk asked to price an option expiring the
day before the event --- an expiry that dodges it --- off the
calendar-interpolated curve quotes essentially the full
\MacClkInterpOverBp\ bp of volatility that does not exist; asked for
one expiring the day after, the one instrument whose whole point is
carrying the event, it quotes nearly the full \MacClkInterpUnderBp\
bp too little.  The errors sit precisely at the maturities an event
trader cares about, they point in opposite directions across one
calendar day, and no quoted expiry ever reveals them --- the rules
agree wherever a price exists.  Interpolation error between quotes
is invisible until it is traded against.

The honest general statement: on a real board the truth between
quotes is not known, and linear-in-$\tau$ is itself a convention ---
the cleanest available, being the constant-accrual statement on the
market's own clock, but a convention nonetheless.  What it
guarantees, and what the calendar rule cannot, is that \emph{known}
events stop leaking into maturities that do not contain them.  The
guarantee is two lines of algebra: between two quotes, $w$ is linear
in $\tau$, and by \eqref{eq:clkclock} $\tau$ grows linearly in $t$
except at an event date, where it jumps by $N_e/365$ --- so,
composed, the dense variance curve is linear in calendar time within
each stretch, pays each event's variance as a jump exactly at its
date, and changes slope only at the quotes.  The remaining question --- where the calendar's dates
and sizes come from --- is the chapter's last construction, and its
most delicate.
